% =============================================================================
% APPENDIX TEMPLATE: Standard Structure for EBF Appendices
% =============================================================================
% 
% This template defines the required structure for all EBF appendices.
% Every appendix should follow this format to ensure consistency.
%
% =============================================================================
%
% IMPORTANT: CORE vs. NON-CORE APPENDICES
% =======================================
%
% CORE appendices (WHO, WHAT, HOW, WHEN) have ADDITIONAL REQUIREMENTS
% beyond the standard template. See Section "CORE APPENDIX REQUIREMENTS"
% at the end of this template.
%
% =============================================================================

% #############################################################################
%
%                           APPENDIX STRUCTURE
%
%    Every EBF Appendix MUST contain the following sections in order:
%
%    ┌─────────────────────────────────────────────────────────────────────┐
%    │  PART 1: FRONT MATTER (Required)                                    │
%    ├─────────────────────────────────────────────────────────────────────┤
%    │  1.1  Header Block (metadata, version, module code)                 │
%    │  1.2  Cross-Reference Map (dependencies to other appendices)        │
%    │  1.3  Chapter Title + Label                                         │
%    │  1.4  Abstract (max 100 words)                                      │
%    │  1.5  Quick Reference Box (key formulas/concepts)                   │
%    │  1.6  Appendix Scope Box (Ziel/In-Scope/Out-of-Scope/Constraints)   │  ← NEW
%    ├─────────────────────────────────────────────────────────────────────┤
%    │  PART 2: CORE CONTENT (Appendix-Specific)                           │
%    ├─────────────────────────────────────────────────────────────────────┤
%    │  2.1  Motivation / The Fundamental Question                         │
%    │  2.2  Core Theory / Main Content                                    │
%    │  2.3  Axioms / Formal Definitions (if applicable)                   │
%    │  2.4  Key Results / Propositions                                    │
%    │  2.5  Integration with Other Appendices                             │
%    ├─────────────────────────────────────────────────────────────────────┤
%    │  PART 3: APPLICATION (Required)                                     │
%    ├─────────────────────────────────────────────────────────────────────┤
%    │  3.1  Worked Example(s)                                             │
%    │  3.2  Implications for Practice                                     │
%    ├─────────────────────────────────────────────────────────────────────┤
%    │  PART 4: BACK MATTER (Required)                                     │
%    ├─────────────────────────────────────────────────────────────────────┤
%    │  4.1  Summary Box                                                   │
%    │  4.2  Glossary of Symbols (+ link to Appendix GLS)           ← MUST   │
%    │  4.3  Critical Foundations / Objections (for CORE)                  │
%    │  4.4  Open Issues / Research Agenda                                 │
%    │  4.5  References (+ link to Master Bibliography)           ← MUST   │
%    │  4.6  Cross-Reference Footer                               ← MUST   │
%    └─────────────────────────────────────────────────────────────────────┘
%
%    MANDATORY LINKS:
%    - Every appendix MUST reference Appendix GLS (Central Glossary)
%    - Every appendix MUST reference bcm2_master_references.bib
%    - Symbols must be consistent with Appendix GLS definitions
%
% #############################################################################


% =============================================================================
% =============================================================================
%
%                    PART 1: FRONT MATTER
%
% =============================================================================
% =============================================================================

% -----------------------------------------------------------------------------
% 1.1 HEADER BLOCK
% -----------------------------------------------------------------------------
% Required metadata for every appendix

% =============================================================================
% APPENDIX [CODE]: [CATEGORY-NAME]: [Descriptive Title]
% =============================================================================
% Module: [BCM2 Module Code] (e.g., BCM2_00_META, BCM2_04_INU_6040)
% Version: [X.Y] ([Month Year])
% Status: [Draft | Review | Final]
%
% [Brief 2-3 line description of what this appendix covers]
%
% =============================================================================

% -----------------------------------------------------------------------------
% 1.2 CROSS-REFERENCE MAP
% -----------------------------------------------------------------------------
% Show dependencies and connections to other appendices/chapters

% =============================================================================
% CROSS-REFERENCE STRUCTURE
% =============================================================================
%
% CRITICAL: Every appendix MUST be linked to its main chapter(s)!
%
% CHAPTER LINKAGE (Required):
% - Primary Chapter: [Chapter X: Title] (the main chapter this appendix supports)
% - Secondary Chapters: [List of other chapters that reference this appendix]
%
% This appendix [provides foundation for / builds upon / integrates with]:
%
% DEPENDENCIES (what this appendix REQUIRES):
% - [Appendix FND]: [What it provides]
% - [Chapter Y]: [What it provides]
%
% DEPENDENTS (what REQUIRES this appendix):
% - [Appendix HAI1]: [How it uses this content]
% - [Chapter W]: [How it uses this content]
%
% ┌────────────────────────────────────────────────────────────────────────────┐
% │ This Appendix    │ Content                    │ Main Chapter Section       │
% ├────────────────────────────────────────────────────────────────────────────┤
% │ Section X.1      │ [Content]                  │ Chapter Y, Section Z       │
% │ Section X.2      │ [Content]                  │ Chapter Y, Section W       │
% └────────────────────────────────────────────────────────────────────────────┘
%
% =============================================================================

% -----------------------------------------------------------------------------
% 1.2a CHAPTER-APPENDIX LINKAGE BOX (Required)
% -----------------------------------------------------------------------------

\begin{tcolorbox}[colback=purple!5!white, colframe=purple!75!black, 
    title=Chapter Linkage]

\textbf{Primary Chapter:} Chapter [X]: [Title]
\begin{itemize}[nosep]
\item Section [X.Y]: [Section title] $\leftrightarrow$ This Appendix Section [A.B]
\item Section [X.Z]: [Section title] $\leftrightarrow$ This Appendix Section [A.C]
\end{itemize}

\textbf{Secondary Chapters:}
\begin{itemize}[nosep]
\item Chapter [M]: [Title] --- references [specific content]
\item Chapter [N]: [Title] --- references [specific content]
\end{itemize}

\textbf{Reading Guide:}
\begin{itemize}[nosep]
\item For conceptual overview: Read Chapter [X] first
\item For formal details: This appendix provides complete specification
\item For proofs: See Appendix [FORMAL-PROOF]
\end{itemize}

\end{tcolorbox}

% -----------------------------------------------------------------------------
% 1.3 CHAPTER TITLE + LABEL
% -----------------------------------------------------------------------------

\chapter{[CATEGORY-NAME]: [Descriptive Title]}
\label{app:[short-label]}


\begin{tcolorbox}[colback=blue!5!white, colframe=blue!75!black,
    title=Quick Reference: Appendix 00]

\textbf{Appendix:} 00 (UNKNOWN)

\textbf{Core Question:} What are the key findings in this domain?

\textbf{Key Concepts:}
\begin{itemize}[nosep]
\item Framework integration
\item Parameter validation
\item Cross-references
\end{itemize}

\textbf{Cross-References:} See related appendices below
\end{tcolorbox}

% -----------------------------------------------------------------------------
% 1.4 ABSTRACT
% -----------------------------------------------------------------------------

\begin{abstract}
[Maximum 100 words describing:]
- What question this appendix answers
- What the main contribution is
- How it connects to EBF
\end{abstract}

% -----------------------------------------------------------------------------
% 1.5 QUICK REFERENCE BOX
% -----------------------------------------------------------------------------

\begin{tcolorbox}[colback=blue!5!white, colframe=blue!75!black, 
    title=Quick Reference: [Appendix Name]]

\textbf{Core Question:} [The fundamental question this appendix answers]

\textbf{Key Formula:}
\begin{equation}
[The most important equation]
\end{equation}

\textbf{Key Concepts:}
\begin{itemize}[nosep]
\item [Concept 1]: [Brief definition]
\item [Concept 2]: [Brief definition]
\item [Concept 3]: [Brief definition]
\end{itemize}

\textbf{Cross-References:} [List of related appendices]
\end{tcolorbox}

% -----------------------------------------------------------------------------
% 1.6 APPENDIX SCOPE BOX (NEW - January 2026)
% -----------------------------------------------------------------------------
% PURPOSE: Define clear boundaries for what this appendix covers and what it
%          delegates to other appendices. This prevents scope creep and ensures
%          clean separation of concerns.
%
% CONSTRAINTS: Appendix-specific requirements and limitations that users must
%              be aware of when applying the content.
%
% MANDATORY FOR: All appendices (METHOD, FORMAL, DOMAIN, CORE, etc.)
% -----------------------------------------------------------------------------

\begin{tcolorbox}[
    colback=orange!5!white,
    colframe=orange!75!black,
    title={\textbf{Appendix Scope: Ziel / In-Scope / Out-of-Scope / Constraints}},
    fonttitle=\bfseries
]

\textbf{Ziel (Objective):}
\begin{itemize}[nosep]
    \item [Primary goal of this appendix in one sentence]
\end{itemize}

\textbf{In-Scope:}
\begin{itemize}[nosep]
    \item [Topic 1 covered in this appendix]
    \item [Topic 2 covered in this appendix]
    \item [Topic 3 covered in this appendix]
\end{itemize}

\textbf{Out-of-Scope (delegiert an andere Appendices):}
\begin{itemize}[nosep]
    \item [Topic A] $\rightarrow$ Appendix [CODE] ([NAME])
    \item [Topic B] $\rightarrow$ Appendix [CODE] ([NAME])
    \item [Topic C] $\rightarrow$ Chapter [N]
\end{itemize}

\textbf{Constraints (Anwendungsgrenzen):}
\begin{itemize}[nosep]
    \item [Constraint 1: When this appendix does NOT apply]
    \item [Constraint 2: Required prerequisites or assumptions]
    \item [Constraint 3: Known limitations of the approach]
\end{itemize}

\textbf{Lieferobjekte (Deliverables):}
\begin{itemize}[nosep]
    \item [Deliverable 1: e.g., Matrix, Table, Formula]
    \item [Deliverable 2: e.g., Algorithm, Axiom System]
    \item [Deliverable 3: e.g., Worked Example, Checklist]
\end{itemize}

\end{tcolorbox}


% =============================================================================
% =============================================================================
%
%                    PART 2: CORE CONTENT
%
% =============================================================================
% =============================================================================

%%%%%%%%%%%%%%%%%%%%%%%%%%%%%%%%%%%%%%%%%%%%%%%%%%%%%%%%%%%%%%%%%%%%%%%%%%%%%%%
\section{The Fundamental Question}
\label{sec:[code]-fundamental}
%%%%%%%%%%%%%%%%%%%%%%%%%%%%%%%%%%%%%%%%%%%%%%%%%%%%%%%%%%%%%%%%%%%%%%%%%%%%%%%

% WHY does this appendix exist? What problem does it solve?

\subsection{Why This Matters}

[Explain the importance of the topic and why EBF needs this appendix]

\subsection{The Core Problem}

[State the specific problem or question this appendix addresses]

\begin{tcolorbox}[colback=red!5!white, colframe=red!75!black, 
    title=The Fundamental Question]
\textbf{[State the question in one clear sentence]}
\end{tcolorbox}


%%%%%%%%%%%%%%%%%%%%%%%%%%%%%%%%%%%%%%%%%%%%%%%%%%%%%%%%%%%%%%%%%%%%%%%%%%%%%%%
\section{Core Theory}
\label{sec:[code]-theory}
%%%%%%%%%%%%%%%%%%%%%%%%%%%%%%%%%%%%%%%%%%%%%%%%%%%%%%%%%%%%%%%%%%%%%%%%%%%%%%%

% The main theoretical content - structure depends on appendix type

\subsection{[Main Concept 1]}

[Detailed explanation with equations, definitions, and discussion]

\subsection{[Main Concept 2]}

[Detailed explanation with equations, definitions, and discussion]


%%%%%%%%%%%%%%%%%%%%%%%%%%%%%%%%%%%%%%%%%%%%%%%%%%%%%%%%%%%%%%%%%%%%%%%%%%%%%%%
\section{Axioms and Formal Definitions}
\label{sec:[code]-axioms}
%%%%%%%%%%%%%%%%%%%%%%%%%%%%%%%%%%%%%%%%%%%%%%%%%%%%%%%%%%%%%%%%%%%%%%%%%%%%%%%

% If applicable - formal axiom system

\begin{tcolorbox}[colback=gray!5!white, colframe=gray!75!black, 
    title=Axiom [CODE]-1: [Name]]
\textbf{Statement:} [Formal statement]

\textbf{Interpretation:} [What it means]

\textbf{Implication:} [What follows from it]
\end{tcolorbox}

% Repeat for each axiom


%%%%%%%%%%%%%%%%%%%%%%%%%%%%%%%%%%%%%%%%%%%%%%%%%%%%%%%%%%%%%%%%%%%%%%%%%%%%%%%
\section{Key Results}
\label{sec:[code]-results}
%%%%%%%%%%%%%%%%%%%%%%%%%%%%%%%%%%%%%%%%%%%%%%%%%%%%%%%%%%%%%%%%%%%%%%%%%%%%%%%

% Main propositions, theorems, or findings

\begin{proposition}[Name]
[Formal statement]
\end{proposition}

\begin{proof}
[Proof or reference to FORMAL-PROOF appendix]
\end{proof}


%%%%%%%%%%%%%%%%%%%%%%%%%%%%%%%%%%%%%%%%%%%%%%%%%%%%%%%%%%%%%%%%%%%%%%%%%%%%%%%
\section{Integration with Other Appendices}
\label{sec:[code]-integration}
%%%%%%%%%%%%%%%%%%%%%%%%%%%%%%%%%%%%%%%%%%%%%%%%%%%%%%%%%%%%%%%%%%%%%%%%%%%%%%%

% How this appendix connects to the broader EBF system

\subsection{Connection to [Related Appendix 1]}

[Explain the relationship, provide formulas showing integration]

\begin{tcolorbox}[colback=yellow!5!white, colframe=yellow!75!black, 
    title=Integration: This Appendix $\leftrightarrow$ [Related Appendix]]
\textbf{What this appendix provides:}
\begin{itemize}[nosep]
\item [Item 1]
\item [Item 2]
\end{itemize}

\textbf{What it requires:}
\begin{itemize}[nosep]
\item [Item 1]
\item [Item 2]
\end{itemize}
\end{tcolorbox}


% =============================================================================
% =============================================================================
%
%                    PART 3: APPLICATION
%
% =============================================================================
% =============================================================================

%%%%%%%%%%%%%%%%%%%%%%%%%%%%%%%%%%%%%%%%%%%%%%%%%%%%%%%%%%%%%%%%%%%%%%%%%%%%%%%
\section{Worked Example}
\label{sec:[code]-example}
%%%%%%%%%%%%%%%%%%%%%%%%%%%%%%%%%%%%%%%%%%%%%%%%%%%%%%%%%%%%%%%%%%%%%%%%%%%%%%%

% At least one detailed worked example showing the theory in action

\subsection{[Example Name]}

\paragraph{Setup.} [Describe the scenario]

\paragraph{Application.} [Show step-by-step how the theory applies]

\paragraph{Result.} [State what we learn]

\paragraph{Discussion.} [Interpret the result]


%%%%%%%%%%%%%%%%%%%%%%%%%%%%%%%%%%%%%%%%%%%%%%%%%%%%%%%%%%%%%%%%%%%%%%%%%%%%%%%
\section{Implications for Practice}
\label{sec:[code]-practice}
%%%%%%%%%%%%%%%%%%%%%%%%%%%%%%%%%%%%%%%%%%%%%%%%%%%%%%%%%%%%%%%%%%%%%%%%%%%%%%%

% Practical takeaways for EBF users

\begin{tcolorbox}[colback=green!5!white, colframe=green!75!black, 
    title=Practical Implications]
\begin{enumerate}
\item \textbf{[Implication 1]:} [Explanation]
\item \textbf{[Implication 2]:} [Explanation]
\item \textbf{[Implication 3]:} [Explanation]
\end{enumerate}
\end{tcolorbox}


% =============================================================================
% =============================================================================
%

%%%%%%%%%%%%%%%%%%%%%%%%%%%%%%%%%%%%%%%%%%%%%%%%%%%%%%%%%%%%%%%%%%%%%%%%%%%%%%%
%%%%%%%%%%%%%%%%%%%%%%%%%%%%%%%%%%%%%%%%%%%%%%%%%%%%%%%%%%%%%%%%%%%%%%%%%%%%%%%
%
%                    PART 4: BACK MATTER (Required)
%
%%%%%%%%%%%%%%%%%%%%%%%%%%%%%%%%%%%%%%%%%%%%%%%%%%%%%%%%%%%%%%%%%%%%%%%%%%%%%%%
%%%%%%%%%%%%%%%%%%%%%%%%%%%%%%%%%%%%%%%%%%%%%%%%%%%%%%%%%%%%%%%%%%%%%%%%%%%%%%%

% =============================================================================
% BACK MATTER REQUIREMENTS
% =============================================================================
%
% Every appendix MUST include:
%   4.1  Summary Box
%   4.2  Glossary of Symbols (appendix-specific + link to central Glossary G)
%   4.3  Critical Foundations / Objections (for CORE appendices)
%   4.4  Open Issues / Research Agenda
%   4.5  References (appendix-specific + link to Master Bibliography)
%
% CRITICAL: All symbol definitions MUST be consistent with Appendix GLS (Glossary)
% CRITICAL: All citations MUST appear in bcm2_master_references.bib
%
% =============================================================================


%%%%%%%%%%%%%%%%%%%%%%%%%%%%%%%%%%%%%%%%%%%%%%%%%%%%%%%%%%%%%%%%%%%%%%%%%%%%%%%
\section{Summary}
\label{sec:[code]-summary}
%%%%%%%%%%%%%%%%%%%%%%%%%%%%%%%%%%%%%%%%%%%%%%%%%%%%%%%%%%%%%%%%%%%%%%%%%%%%%%%

\begin{tcolorbox}[colback=green!10!white, colframe=green!75!black, 
    title=\textbf{Appendix [CODE]: Summary}]

\textbf{Core Question:} [What question does this appendix answer?]

\textbf{Main Contribution:}
\begin{itemize}[nosep]
    \item [Key result 1]
    \item [Key result 2]
    \item [Key result 3]
\end{itemize}

\textbf{Key Outputs:}
\begin{itemize}[nosep]
    \item [Mathematical object 1]: [Brief description]
    \item [Mathematical object 2]: [Brief description]
\end{itemize}

\textbf{Integration:} This appendix connects to [list CORE appendices] 
via [describe connection].

\end{tcolorbox}


%%%%%%%%%%%%%%%%%%%%%%%%%%%%%%%%%%%%%%%%%%%%%%%%%%%%%%%%%%%%%%%%%%%%%%%%%%%%%%%
\section{Glossary of Symbols}
\label{sec:[code]-glossary}
%%%%%%%%%%%%%%%%%%%%%%%%%%%%%%%%%%%%%%%%%%%%%%%%%%%%%%%%%%%%%%%%%%%%%%%%%%%%%%%

% -----------------------------------------------------------------------------
% INSTRUCTION: Include ONLY symbols introduced or modified in THIS appendix.
% All standard EBF symbols are defined in Appendix GLS (Central Glossary).
% 
% If you introduce a NEW symbol:
%   1. Define it here
%   2. Ensure it is added to Appendix GLS in the next update
%
% If you use a STANDARD symbol with a specialized meaning:
%   1. Define the specialization here
%   2. Reference the general definition in Appendix GLS
% -----------------------------------------------------------------------------

\begin{tcolorbox}[colback=blue!5!white, colframe=blue!75!black]
\textbf{Central Reference:} For complete symbol definitions, see 
\textbf{Appendix GLS} (\textit{Glossary and Notation}, \S\ref{app:glossary}).

This section lists only symbols introduced or specialized in this appendix.
\end{tcolorbox}

\vspace{0.5em}

\begin{table}[htbp]
\centering
\caption{Symbols Introduced in Appendix [CODE]}
\label{tab:[code]-symbols}
\renewcommand{\arraystretch}{1.3}
\begin{tabular}{@{}clp{6cm}c@{}}
\toprule
\textbf{Symbol} & \textbf{Name} & \textbf{Definition} & \textbf{Also in G?} \\
\midrule
$[symbol]$ & [Name] & [Definition] & \checkmark / NEW \\
$[symbol]$ & [Name] & [Definition] & \checkmark / NEW \\
\bottomrule
\end{tabular}
\end{table}

% Mark symbols as:
%   \checkmark = Standard symbol defined in Appendix GLS (reference only)
%   NEW = Novel symbol introduced here (must be added to G)


%%%%%%%%%%%%%%%%%%%%%%%%%%%%%%%%%%%%%%%%%%%%%%%%%%%%%%%%%%%%%%%%%%%%%%%%%%%%%%%
\section{Critical Foundations and Objections}
\label{sec:[code]-foundations}
%%%%%%%%%%%%%%%%%%%%%%%%%%%%%%%%%%%%%%%%%%%%%%%%%%%%%%%%%%%%%%%%%%%%%%%%%%%%%%%

% Required for CORE appendices; optional for others

\subsection{Objection 1: [Name]}

\textbf{Objection:} [State the objection]

\textbf{Response:} [Provide the response]

\subsection{Objection 2: [Name]}

\textbf{Objection:} [State the objection]

\textbf{Response:} [Provide the response]


%%%%%%%%%%%%%%%%%%%%%%%%%%%%%%%%%%%%%%%%%%%%%%%%%%%%%%%%%%%%%%%%%%%%%%%%%%%%%%%
\section{Open Issues and Research Agenda}
\label{sec:[code]-open}
%%%%%%%%%%%%%%%%%%%%%%%%%%%%%%%%%%%%%%%%%%%%%%%%%%%%%%%%%%%%%%%%%%%%%%%%%%%%%%%

\subsection{Methodological Priorities}

\begin{enumerate}
    \item \textbf{[Issue 1]:} [Description]
    \item \textbf{[Issue 2]:} [Description]
\end{enumerate}

\subsection{Empirical Priorities}

\begin{enumerate}
    \item \textbf{[Issue 1]:} [Description]
    \item \textbf{[Issue 2]:} [Description]
\end{enumerate}


%%%%%%%%%%%%%%%%%%%%%%%%%%%%%%%%%%%%%%%%%%%%%%%%%%%%%%%%%%%%%%%%%%%%%%%%%%%%%%%
\section{References}
\label{sec:[code]-references}
%%%%%%%%%%%%%%%%%%%%%%%%%%%%%%%%%%%%%%%%%%%%%%%%%%%%%%%%%%%%%%%%%%%%%%%%%%%%%%%

% -----------------------------------------------------------------------------
% INSTRUCTION: Reference structure for EBF appendices
%
% 1. MASTER BIBLIOGRAPHY: bcm2_master_references.bib
%    Contains all shared references across the framework.
%    Use \cite{key} for standard citations.
%
% 2. APPENDIX-SPECIFIC: References unique to this appendix
%    Include here only if not in master bibliography.
%
% 3. DATA SOURCES: Empirical data used in this appendix
%    List with URLs and access dates.
%
% CONSISTENCY CHECK: Before finalizing, verify all citations appear in
% either the master bibliography or this section.
% -----------------------------------------------------------------------------

\begin{tcolorbox}[colback=blue!5!white, colframe=blue!75!black]
\textbf{Master Bibliography:} For complete reference details, see 
\texttt{bcm2\_master\_references.bib} in the repository.

This section lists appendix-specific references and data sources.
\end{tcolorbox}

\vspace{0.5em}

\subsection{Key References for This Appendix}

The following works are particularly relevant for the content of this appendix:

\begin{itemize}[nosep]
    \item \textbf{[Topic 1]:} \citet{key1}, \citet{key2}
    \item \textbf{[Topic 2]:} \citet{key3}
\end{itemize}

\subsection{Appendix-Specific References}

% References NOT in the master bibliography

\begin{thebibliography}{99}

\bibitem{[key]} [Author] ([Year]). [Title]. \textit{[Journal]}, [Volume], [Pages].

\end{thebibliography}

\subsection{Data Sources}

\begin{itemize}[nosep]
    \item \textbf{[Data Name]:} [URL] (accessed [Date])
\end{itemize}


%%%%%%%%%%%%%%%%%%%%%%%%%%%%%%%%%%%%%%%%%%%%%%%%%%%%%%%%%%%%%%%%%%%%%%%%%%%%%%%
% CROSS-REFERENCE FOOTER (Required)
%%%%%%%%%%%%%%%%%%%%%%%%%%%%%%%%%%%%%%%%%%%%%%%%%%%%%%%%%%%%%%%%%%%%%%%%%%%%%%%

\vspace{1em}
\hrule
\vspace{0.5em}

\noindent\textit{Cross-references:} 
Appendix GLS (Central Glossary), 
[Appendix FND] ([Topic]), 
[Appendix CIA] ([Topic]).

\noindent\textit{Master Bibliography:} \texttt{bcm2\_master\_references.bib}


% =============================================================================
% END OF APPENDIX [CODE]
% =============================================================================



CORE APPENDIX REQUIREMENTS"
% at the end of this template.
%
% =============================================================================

% #############################################################################
%
%                           APPENDIX STRUCTURE
%
%    Every EBF Appendix MUST contain the following sections in order:
%
%    ┌─────────────────────────────────────────────────────────────────────┐
%    │  PART 1: FRONT MATTER (Required)                                    │
%    ├─────────────────────────────────────────────────────────────────────┤
%    │  1.1  Header Block (metadata, version, module code)                 │
%    │  1.2  Cross-Reference Map (dependencies to other appendices)        │
%    │  1.3  Chapter Title + Label                                         │
%    │  1.4  Abstract (max 100 words)                                      │
%    │  1.5  Quick Reference Box (key formulas/concepts)                   │
%    │  1.6  Appendix Scope Box (Ziel/In-Scope/Out-of-Scope/Constraints)   │  ← NEW
%    ├─────────────────────────────────────────────────────────────────────┤
%    │  PART 2: CORE CONTENT (Appendix-Specific)                           │
%    ├─────────────────────────────────────────────────────────────────────┤
%    │  2.1  Motivation / The Fundamental Question                         │
%    │  2.2  Core Theory / Main Content                                    │
%    │  2.3  Axioms / Formal Definitions (if applicable)                   │
%    │  2.4  Key Results / Propositions                                    │
%    │  2.5  Integration with Other Appendices                             │
%    ├─────────────────────────────────────────────────────────────────────┤
%    │  PART 3: APPLICATION (Required)                                     │
%    ├─────────────────────────────────────────────────────────────────────┤
%    │  3.1  Worked Example(s)                                             │
%    │  3.2  Implications for Practice                                     │
%    ├─────────────────────────────────────────────────────────────────────┤
%    │  PART 4: BACK MATTER (Required)                                     │
%    ├─────────────────────────────────────────────────────────────────────┤
%    │  4.1  Summary Box                                                   │
%    │  4.2  Glossary of Symbols (+ link to Appendix GLS)           ← MUST   │
%    │  4.3  Critical Foundations / Objections (for CORE)                  │
%    │  4.4  Open Issues / Research Agenda                                 │
%    │  4.5  References (+ link to Master Bibliography)           ← MUST   │
%    │  4.6  Cross-Reference Footer                               ← MUST   │
%    └─────────────────────────────────────────────────────────────────────┘
%
%    MANDATORY LINKS:
%    - Every appendix MUST reference Appendix GLS (Central Glossary)
%    - Every appendix MUST reference bcm2_master_references.bib
%    - Symbols must be consistent with Appendix GLS definitions
%
% #############################################################################


% =============================================================================
% =============================================================================
%
%                    PART 1: FRONT MATTER
%
% =============================================================================
% =============================================================================

% -----------------------------------------------------------------------------
% 1.1 HEADER BLOCK
% -----------------------------------------------------------------------------
% Required metadata for every appendix

% =============================================================================
% APPENDIX [CODE]: [CATEGORY-NAME]: [Descriptive Title]
% =============================================================================
% Module: [BCM2 Module Code] (e.g., BCM2_00_META, BCM2_04_INU_6040)
% Version: [X.Y] ([Month Year])
% Status: [Draft | Review | Final]
%
% [Brief 2-3 line description of what this appendix covers]
%
% =============================================================================

% -----------------------------------------------------------------------------
% 1.2 CROSS-REFERENCE MAP
% -----------------------------------------------------------------------------
% Show dependencies and connections to other appendices/chapters

% =============================================================================
% CROSS-REFERENCE STRUCTURE
% =============================================================================
%
% CRITICAL: Every appendix MUST be linked to its main chapter(s)!
%
% CHAPTER LINKAGE (Required):
% - Primary Chapter: [Chapter X: Title] (the main chapter this appendix supports)
% - Secondary Chapters: [List of other chapters that reference this appendix]
%
% This appendix [provides foundation for / builds upon / integrates with]:
%
% DEPENDENCIES (what this appendix REQUIRES):
% - [Appendix FND]: [What it provides]
% - [Chapter Y]: [What it provides]
%
% DEPENDENTS (what REQUIRES this appendix):
% - [Appendix HAI1]: [How it uses this content]
% - [Chapter W]: [How it uses this content]
%
% ┌────────────────────────────────────────────────────────────────────────────┐
% │ This Appendix    │ Content                    │ Main Chapter Section       │
% ├────────────────────────────────────────────────────────────────────────────┤
% │ Section X.1      │ [Content]                  │ Chapter Y, Section Z       │
% │ Section X.2      │ [Content]                  │ Chapter Y, Section W       │
% └────────────────────────────────────────────────────────────────────────────┘
%
% =============================================================================

% -----------------------------------------------------------------------------
% 1.2a CHAPTER-APPENDIX LINKAGE BOX (Required)
% -----------------------------------------------------------------------------

\begin{tcolorbox}[colback=purple!5!white, colframe=purple!75!black, 
    title=Chapter Linkage]

\textbf{Primary Chapter:} Chapter [X]: [Title]
\begin{itemize}[nosep]
\item Section [X.Y]: [Section title] $\leftrightarrow$ This Appendix Section [A.B]
\item Section [X.Z]: [Section title] $\leftrightarrow$ This Appendix Section [A.C]
\end{itemize}

\textbf{Secondary Chapters:}
\begin{itemize}[nosep]
\item Chapter [M]: [Title] --- references [specific content]
\item Chapter [N]: [Title] --- references [specific content]
\end{itemize}

\textbf{Reading Guide:}
\begin{itemize}[nosep]
\item For conceptual overview: Read Chapter [X] first
\item For formal details: This appendix provides complete specification
\item For proofs: See Appendix [FORMAL-PROOF]
\end{itemize}

\end{tcolorbox}

% -----------------------------------------------------------------------------
% 1.3 CHAPTER TITLE + LABEL
% -----------------------------------------------------------------------------

\chapter{[CATEGORY-NAME]: [Descriptive Title]}
\label{app:[short-label]}

% -----------------------------------------------------------------------------
% 1.4 ABSTRACT
% -----------------------------------------------------------------------------

\begin{abstract}
[Maximum 100 words describing:]
- What question this appendix answers
- What the main contribution is
- How it connects to EBF
\end{abstract}

% -----------------------------------------------------------------------------
% 1.5 QUICK REFERENCE BOX
% -----------------------------------------------------------------------------

\begin{tcolorbox}[colback=blue!5!white, colframe=blue!75!black, 
    title=Quick Reference: [Appendix Name]]

\textbf{Core Question:} [The fundamental question this appendix answers]

\textbf{Key Formula:}
\begin{equation}
[The most important equation]
\end{equation}

\textbf{Key Concepts:}
\begin{itemize}[nosep]
\item [Concept 1]: [Brief definition]
\item [Concept 2]: [Brief definition]
\item [Concept 3]: [Brief definition]
\end{itemize}

\textbf{Cross-References:} [List of related appendices]
\end{tcolorbox}

% -----------------------------------------------------------------------------
% 1.6 APPENDIX SCOPE BOX (NEW - January 2026)
% -----------------------------------------------------------------------------
% PURPOSE: Define clear boundaries for what this appendix covers and what it
%          delegates to other appendices. This prevents scope creep and ensures
%          clean separation of concerns.
%
% CONSTRAINTS: Appendix-specific requirements and limitations that users must
%              be aware of when applying the content.
%
% MANDATORY FOR: All appendices (METHOD, FORMAL, DOMAIN, CORE, etc.)
% -----------------------------------------------------------------------------

\begin{tcolorbox}[
    colback=orange!5!white,
    colframe=orange!75!black,
    title={\textbf{Appendix Scope: Ziel / In-Scope / Out-of-Scope / Constraints}},
    fonttitle=\bfseries
]

\textbf{Ziel (Objective):}
\begin{itemize}[nosep]
    \item [Primary goal of this appendix in one sentence]
\end{itemize}

\textbf{In-Scope:}
\begin{itemize}[nosep]
    \item [Topic 1 covered in this appendix]
    \item [Topic 2 covered in this appendix]
    \item [Topic 3 covered in this appendix]
\end{itemize}

\textbf{Out-of-Scope (delegiert an andere Appendices):}
\begin{itemize}[nosep]
    \item [Topic A] $\rightarrow$ Appendix [CODE] ([NAME])
    \item [Topic B] $\rightarrow$ Appendix [CODE] ([NAME])
    \item [Topic C] $\rightarrow$ Chapter [N]
\end{itemize}

\textbf{Constraints (Anwendungsgrenzen):}
\begin{itemize}[nosep]
    \item [Constraint 1: When this appendix does NOT apply]
    \item [Constraint 2: Required prerequisites or assumptions]
    \item [Constraint 3: Known limitations of the approach]
\end{itemize}

\textbf{Lieferobjekte (Deliverables):}
\begin{itemize}[nosep]
    \item [Deliverable 1: e.g., Matrix, Table, Formula]
    \item [Deliverable 2: e.g., Algorithm, Axiom System]
    \item [Deliverable 3: e.g., Worked Example, Checklist]
\end{itemize}

\end{tcolorbox}


% =============================================================================
% =============================================================================
%
%                    PART 2: CORE CONTENT
%
% =============================================================================
% =============================================================================

%%%%%%%%%%%%%%%%%%%%%%%%%%%%%%%%%%%%%%%%%%%%%%%%%%%%%%%%%%%%%%%%%%%%%%%%%%%%%%%
\section{The Fundamental Question}
\label{sec:[code]-fundamental}
%%%%%%%%%%%%%%%%%%%%%%%%%%%%%%%%%%%%%%%%%%%%%%%%%%%%%%%%%%%%%%%%%%%%%%%%%%%%%%%

% WHY does this appendix exist? What problem does it solve?

\subsection{Why This Matters}

[Explain the importance of the topic and why EBF needs this appendix]

\subsection{The Core Problem}

[State the specific problem or question this appendix addresses]

\begin{tcolorbox}[colback=red!5!white, colframe=red!75!black, 
    title=The Fundamental Question]
\textbf{[State the question in one clear sentence]}
\end{tcolorbox}


%%%%%%%%%%%%%%%%%%%%%%%%%%%%%%%%%%%%%%%%%%%%%%%%%%%%%%%%%%%%%%%%%%%%%%%%%%%%%%%
\section{Core Theory}
\label{sec:[code]-theory}
%%%%%%%%%%%%%%%%%%%%%%%%%%%%%%%%%%%%%%%%%%%%%%%%%%%%%%%%%%%%%%%%%%%%%%%%%%%%%%%

% The main theoretical content - structure depends on appendix type

\subsection{[Main Concept 1]}

[Detailed explanation with equations, definitions, and discussion]

\subsection{[Main Concept 2]}

[Detailed explanation with equations, definitions, and discussion]


%%%%%%%%%%%%%%%%%%%%%%%%%%%%%%%%%%%%%%%%%%%%%%%%%%%%%%%%%%%%%%%%%%%%%%%%%%%%%%%
\section{Axioms and Formal Definitions}
\label{sec:[code]-axioms}
%%%%%%%%%%%%%%%%%%%%%%%%%%%%%%%%%%%%%%%%%%%%%%%%%%%%%%%%%%%%%%%%%%%%%%%%%%%%%%%

% If applicable - formal axiom system

\begin{tcolorbox}[colback=gray!5!white, colframe=gray!75!black, 
    title=Axiom [CODE]-1: [Name]]
\textbf{Statement:} [Formal statement]

\textbf{Interpretation:} [What it means]

\textbf{Implication:} [What follows from it]
\end{tcolorbox}

% Repeat for each axiom


%%%%%%%%%%%%%%%%%%%%%%%%%%%%%%%%%%%%%%%%%%%%%%%%%%%%%%%%%%%%%%%%%%%%%%%%%%%%%%%
\section{Key Results}
\label{sec:[code]-results}
%%%%%%%%%%%%%%%%%%%%%%%%%%%%%%%%%%%%%%%%%%%%%%%%%%%%%%%%%%%%%%%%%%%%%%%%%%%%%%%

% Main propositions, theorems, or findings

\begin{proposition}[Name]
[Formal statement]
\end{proposition}

\begin{proof}
[Proof or reference to FORMAL-PROOF appendix]
\end{proof}


%%%%%%%%%%%%%%%%%%%%%%%%%%%%%%%%%%%%%%%%%%%%%%%%%%%%%%%%%%%%%%%%%%%%%%%%%%%%%%%
\section{Integration with Other Appendices}
\label{sec:[code]-integration}
%%%%%%%%%%%%%%%%%%%%%%%%%%%%%%%%%%%%%%%%%%%%%%%%%%%%%%%%%%%%%%%%%%%%%%%%%%%%%%%

% How this appendix connects to the broader EBF system

\subsection{Connection to [Related Appendix 1]}

[Explain the relationship, provide formulas showing integration]

\begin{tcolorbox}[colback=yellow!5!white, colframe=yellow!75!black, 
    title=Integration: This Appendix $\leftrightarrow$ [Related Appendix]]
\textbf{What this appendix provides:}
\begin{itemize}[nosep]
\item [Item 1]
\item [Item 2]
\end{itemize}

\textbf{What it requires:}
\begin{itemize}[nosep]
\item [Item 1]
\item [Item 2]
\end{itemize}
\end{tcolorbox}


% =============================================================================
% =============================================================================
%
%                    PART 3: APPLICATION
%
% =============================================================================
% =============================================================================

%%%%%%%%%%%%%%%%%%%%%%%%%%%%%%%%%%%%%%%%%%%%%%%%%%%%%%%%%%%%%%%%%%%%%%%%%%%%%%%
\section{Worked Example}
\label{sec:[code]-example}
%%%%%%%%%%%%%%%%%%%%%%%%%%%%%%%%%%%%%%%%%%%%%%%%%%%%%%%%%%%%%%%%%%%%%%%%%%%%%%%

% At least one detailed worked example showing the theory in action

\subsection{[Example Name]}

\paragraph{Setup.} [Describe the scenario]

\paragraph{Application.} [Show step-by-step how the theory applies]

\paragraph{Result.} [State what we learn]

\paragraph{Discussion.} [Interpret the result]


%%%%%%%%%%%%%%%%%%%%%%%%%%%%%%%%%%%%%%%%%%%%%%%%%%%%%%%%%%%%%%%%%%%%%%%%%%%%%%%
\section{Implications for Practice}
\label{sec:[code]-practice}
%%%%%%%%%%%%%%%%%%%%%%%%%%%%%%%%%%%%%%%%%%%%%%%%%%%%%%%%%%%%%%%%%%%%%%%%%%%%%%%

% Practical takeaways for EBF users

\begin{tcolorbox}[colback=green!5!white, colframe=green!75!black, 
    title=Practical Implications]
\begin{enumerate}
\item \textbf{[Implication 1]:} [Explanation]
\item \textbf{[Implication 2]:} [Explanation]
\item \textbf{[Implication 3]:} [Explanation]
\end{enumerate}
\end{tcolorbox}


% =============================================================================
% =============================================================================
%
%                    PART 4: BACK MATTER
%
% =============================================================================
% =============================================================================

%%%%%%%%%%%%%%%%%%%%%%%%%%%%%%%%%%%%%%%%%%%%%%%%%%%%%%%%%%%%%%%%%%%%%%%%%%%%%%%
\section{Summary}
\label{sec:[code]-summary}
%%%%%%%%%%%%%%%%%%%%%%%%%%%%%%%%%%%%%%%%%%%%%%%%%%%%%%%%%%%%%%%%%%%%%%%%%%%%%%%

\begin{tcolorbox}[colback=gray!5!white, colframe=gray!75!black, 
    title=Appendix [CODE] Summary: [Name]]

\textbf{Core Contribution:}
[One paragraph summary]

\textbf{Key Formula:}
\begin{equation}
[Most important equation]
\end{equation}

\textbf{Key Insights:}
\begin{itemize}[nosep]
\item [Insight 1]
\item [Insight 2]
\item [Insight 3]
\end{itemize}

\textbf{Next Steps:} See [Related Appendix] for [continuation]
\end{tcolorbox}


%%%%%%%%%%%%%%%%%%%%%%%%%%%%%%%%%%%%%%%%%%%%%%%%%%%%%%%%%%%%%%%%%%%%%%%%%%%%%%%
\section{Glossary of Symbols}
\label{sec:[code]-glossary}
%%%%%%%%%%%%%%%%%%%%%%%%%%%%%%%%%%%%%%%%%%%%%%%%%%%%%%%%%%%%%%%%%%%%%%%%%%%%%%%

\begin{table}[htbp]
\centering
\caption{Symbols in [Appendix Name]}
\label{tab:[code]-symbols}
\renewcommand{\arraystretch}{1.3}
\begin{tabular}{@{}lll@{}}
\toprule
\textbf{Symbol} & \textbf{Name} & \textbf{Definition} \\
\midrule
$[symbol]$ & [Name] & [Definition] \\
$[symbol]$ & [Name] & [Definition] \\
\bottomrule
\end{tabular}
\end{table}


%%%%%%%%%%%%%%%%%%%%%%%%%%%%%%%%%%%%%%%%%%%%%%%%%%%%%%%%%%%%%%%%%%%%%%%%%%%%%%%
\section{Critical Foundations and Objections}
\label{sec:[code]-critical}
%%%%%%%%%%%%%%%%%%%%%%%%%%%%%%%%%%%%%%%%%%%%%%%%%%%%%%%%%%%%%%%%%%%%%%%%%%%%%%%

% For CORE appendices: Address potential objections

\subsection{Objection 1: [Name]}

\begin{tcolorbox}[colback=red!5!white, colframe=red!50!black, 
    title=Objection: [Strongest form]]
[State the objection in its strongest form]
\end{tcolorbox}

\textbf{Response:} [Provide rigorous response with evidence]


%%%%%%%%%%%%%%%%%%%%%%%%%%%%%%%%%%%%%%%%%%%%%%%%%%%%%%%%%%%%%%%%%%%%%%%%%%%%%%%
\section{Open Issues and Research Agenda}
\label{sec:[code]-open}
%%%%%%%%%%%%%%%%%%%%%%%%%%%%%%%%%%%%%%%%%%%%%%%%%%%%%%%%%%%%%%%%%%%%%%%%%%%%%%%

% Honest acknowledgment of what remains to be done

\begin{table}[htbp]
\centering
\caption{Research Roadmap for [Appendix Name]}
\label{tab:[code]-roadmap}
\renewcommand{\arraystretch}{1.3}
\begin{tabular}{@{}clcl@{}}
\toprule
\textbf{\#} & \textbf{Issue} & \textbf{Priority} & \textbf{Status} \\
\midrule
1 & [Issue 1] & HIGH & [Status] \\
2 & [Issue 2] & MEDIUM & [Status] \\
3 & [Issue 3] & LOW & [Status] \\
\bottomrule
\end{tabular}
\end{table}


%%%%%%%%%%%%%%%%%%%%%%%%%%%%%%%%%%%%%%%%%%%%%%%%%%%%%%%%%%%%%%%%%%%%%%%%%%%%%%%
\section{References}
\label{sec:[code]-references}
%%%%%%%%%%%%%%%%%%%%%%%%%%%%%%%%%%%%%%%%%%%%%%%%%%%%%%%%%%%%%%%%%%%%%%%%%%%%%%%

% Appendix-specific references

\begin{thebibliography}{99}

\bibitem{[key]} [Author] ([Year]). [Title]. \textit{[Journal]}, [Volume], [Pages].

\end{thebibliography}


% =============================================================================
% END OF APPENDIX [CODE]
% =============================================================================


% =============================================================================
% =============================================================================
%
%                    CORE APPENDIX REQUIREMENTS
%
%              Additional requirements for CORE-WHO, CORE-WHAT,
%                       CORE-HOW, and CORE-WHEN
%
% =============================================================================
% =============================================================================

% CORE appendices are FUNDAMENTAL to EBF - without them the framework
% does not function. They must meet ALL standard requirements PLUS these
% additional CORE-specific requirements:

%%%%%%%%%%%%%%%%%%%%%%%%%%%%%%%%%%%%%%%%%%%%%%%%%%%%%%%%%%%%%%%%%%%%%%%%%%%%%%%
% CORE REQUIREMENT 1: ANSWERS A FUNDAMENTAL QUESTION
%%%%%%%%%%%%%%%%%%%%%%%%%%%%%%%%%%%%%%%%%%%%%%%%%%%%%%%%%%%%%%%%%%%%%%%%%%%%%%%

% Each CORE appendix must answer exactly ONE of the four fundamental questions:
%
% ┌─────────────┬────────────────────────────┬──────────────────────────────┐
% │ CORE        │ Fundamental Question       │ What it defines              │
% ├─────────────┼────────────────────────────┼──────────────────────────────┤
% │ CORE-WHO    │ Who can have utility?      │ Levels L = 0, 1, ..., ∞      │
% │ CORE-WHAT   │ What constitutes utility?  │ Dimensions d ∈ FEPSDE        │
% │ CORE-HOW    │ How do components interact?│ Complementarity γ            │
% │ CORE-WHEN   │ When does context matter?  │ Context Ψ = (Ψ₁, ..., Ψ₈)   │
% └─────────────┴────────────────────────────┴──────────────────────────────┘

%%%%%%%%%%%%%%%%%%%%%%%%%%%%%%%%%%%%%%%%%%%%%%%%%%%%%%%%%%%%%%%%%%%%%%%%%%%%%%%
% CORE REQUIREMENT 2: COMPLETE AXIOM SYSTEM
%%%%%%%%%%%%%%%%%%%%%%%%%%%%%%%%%%%%%%%%%%%%%%%%%%%%%%%%%%%%%%%%%%%%%%%%%%%%%%%

% Every CORE appendix MUST define its own axiom system with:
%   - Unique axiom codes (e.g., SRL-1, γ-1, Ψ-1)
%   - Formal mathematical statements
%   - Interpretations and implications
%   - Cross-references to related axioms in other CORE appendices
%
% Minimum axiom counts:
%   - CORE-WHO:  8+ axioms (SRL-1 to SRL-8)
%   - CORE-WHAT: 200+ axioms (INU, KNU, IDN combined)
%   - CORE-HOW:  10+ axioms (γ-1 to γ-10)
%   - CORE-WHEN: 8+ axioms (Ψ-1 to Ψ-8)

\begin{tcolorbox}[colback=red!5!white, colframe=red!75!black, 
    title=CORE Axiom Template]
    
\textbf{Axiom [PREFIX]-[NUMBER]: [Name]}

\textbf{Statement:}
\begin{equation}
[Formal mathematical statement]
\end{equation}

\textbf{Interpretation:} [What this axiom means in plain language]

\textbf{Implication:} [What follows from this axiom]

\textbf{Testable Prediction:} [How this axiom could be falsified]

\textbf{Cross-References:} [Related axioms in other CORE appendices]

\end{tcolorbox}

%%%%%%%%%%%%%%%%%%%%%%%%%%%%%%%%%%%%%%%%%%%%%%%%%%%%%%%%%%%%%%%%%%%%%%%%%%%%%%%
% CORE REQUIREMENT 3: COMPREHENSIVE LITERATURE REVIEW
%%%%%%%%%%%%%%%%%%%%%%%%%%%%%%%%%%%%%%%%%%%%%%%%%%%%%%%%%%%%%%%%%%%%%%%%%%%%%%%

% CORE appendices must provide EXHAUSTIVE literature coverage organized by:
%
% 1. FOUNDATIONAL PAPERS (Classics)
%    - The seminal works that established the field
%    - Nobel Prize winning research where applicable
%
% 2. RECENT ADVANCES (2020-2025)
%    - State of the art research
%    - New empirical findings
%
% 3. EMPIRICAL EVIDENCE
%    - Laboratory experiments
%    - Field studies
%    - Natural experiments
%
% 4. CRITIQUES AND DEBATES
%    - Major objections in the literature
%    - Ongoing controversies
%
% 5. CROSS-DISCIPLINARY SOURCES
%    - Psychology, sociology, neuroscience
%    - Philosophy, political science
%
% Minimum reference counts:
%   - CORE-WHO:  100+ references
%   - CORE-WHAT: 150+ references
%   - CORE-HOW:  80+ references
%   - CORE-WHEN: 80+ references

%%%%%%%%%%%%%%%%%%%%%%%%%%%%%%%%%%%%%%%%%%%%%%%%%%%%%%%%%%%%%%%%%%%%%%%%%%%%%%%
% CORE REQUIREMENT 4: CRITICAL FOUNDATIONS (Objections & Responses)
%%%%%%%%%%%%%%%%%%%%%%%%%%%%%%%%%%%%%%%%%%%%%%%%%%%%%%%%%%%%%%%%%%%%%%%%%%%%%%%

% CORE appendices must anticipate and respond to major objections.
% Minimum: 5-10 objections, each with:

\begin{tcolorbox}[colback=yellow!5!white, colframe=yellow!75!black, 
    title=Critical Foundation Template]

\textbf{Objection [N]: [Name of objection]}

\textbf{Source:} [Who makes this objection - discipline/author]

\textbf{Strongest Form:}
\begin{quote}
[State the objection in its most powerful form - steelman it]
\end{quote}

\textbf{Response:}
[Provide rigorous, evidence-based response]

\textbf{Empirical Evidence:}
[Cite studies that address this objection]

\textbf{Remaining Limitations:}
[Honest acknowledgment of what the objection still challenges]

\end{tcolorbox}

%%%%%%%%%%%%%%%%%%%%%%%%%%%%%%%%%%%%%%%%%%%%%%%%%%%%%%%%%%%%%%%%%%%%%%%%%%%%%%%
% CORE REQUIREMENT 5: OPEN ISSUES AND RESEARCH AGENDA
%%%%%%%%%%%%%%%%%%%%%%%%%%%%%%%%%%%%%%%%%%%%%%%%%%%%%%%%%%%%%%%%%%%%%%%%%%%%%%%

% CORE appendices must honestly acknowledge gaps and provide research roadmap.
% Minimum: 5-8 open issues with:

\begin{table}[htbp]
\centering
\caption{Research Roadmap Template for CORE Appendix}
\renewcommand{\arraystretch}{1.3}
\begin{tabular}{@{}clccl@{}}
\toprule
\textbf{\#} & \textbf{Issue} & \textbf{Priority} & \textbf{Status} & \textbf{Requirements} \\
\midrule
1 & [Issue name] & 🔴 HIGH & [Status] & [What's needed to resolve] \\
2 & [Issue name] & 🟡 MEDIUM & [Status] & [What's needed to resolve] \\
3 & [Issue name] & 🟢 LOW & [Status] & [What's needed to resolve] \\
\bottomrule
\end{tabular}
\end{table}

%%%%%%%%%%%%%%%%%%%%%%%%%%%%%%%%%%%%%%%%%%%%%%%%%%%%%%%%%%%%%%%%%%%%%%%%%%%%%%%
% CORE REQUIREMENT 6: BIDIRECTIONAL INTEGRATION WITH ALL OTHER CORE
%%%%%%%%%%%%%%%%%%%%%%%%%%%%%%%%%%%%%%%%%%%%%%%%%%%%%%%%%%%%%%%%%%%%%%%%%%%%%%%

% Every CORE appendix must formally connect to ALL other CORE appendices:
%
% Required integration sections:
%   - CORE-WHO must connect to: WHAT, HOW, WHEN
%   - CORE-WHAT must connect to: WHO, HOW, WHEN
%   - CORE-HOW must connect to: WHO, WHAT, WHEN
%   - CORE-WHEN must connect to: WHO, WHAT, HOW
%
% Each integration must specify:
%   1. What this CORE provides to the other
%   2. What this CORE requires from the other
%   3. The mathematical interface (shared variables, constraints)
%   4. Consistency requirements

\begin{tcolorbox}[colback=blue!5!white, colframe=blue!75!black, 
    title=CORE Integration Template: [This CORE] $\leftrightarrow$ [Other CORE]]

\textbf{What [This CORE] provides:}
\begin{itemize}[nosep]
\item [Variable/concept 1]
\item [Variable/concept 2]
\end{itemize}

\textbf{What [This CORE] requires:}
\begin{itemize}[nosep]
\item [Variable/concept 1]
\item [Variable/concept 2]
\end{itemize}

\textbf{Mathematical Interface:}
\begin{equation}
[Equation showing how variables connect]
\end{equation}

\textbf{Consistency Constraint:}
\begin{equation}
[Constraint that must hold for consistency]
\end{equation}

\end{tcolorbox}

%%%%%%%%%%%%%%%%%%%%%%%%%%%%%%%%%%%%%%%%%%%%%%%%%%%%%%%%%%%%%%%%%%%%%%%%%%%%%%%
% CORE REQUIREMENT 7: COMPREHENSIVE GLOSSARY
%%%%%%%%%%%%%%%%%%%%%%%%%%%%%%%%%%%%%%%%%%%%%%%%%%%%%%%%%%%%%%%%%%%%%%%%%%%%%%%

% CORE appendices must define ALL symbols and terms:
%   - Every Greek letter used
%   - Every subscript/superscript notation
%   - Every technical term
%   - Cross-references to where each is defined
%
% Minimum: 20+ symbol definitions

%%%%%%%%%%%%%%%%%%%%%%%%%%%%%%%%%%%%%%%%%%%%%%%%%%%%%%%%%%%%%%%%%%%%%%%%%%%%%%%
% SUMMARY: CORE APPENDIX CHECKLIST
%%%%%%%%%%%%%%%%%%%%%%%%%%%%%%%%%%%%%%%%%%%%%%%%%%%%%%%%%%%%%%%%%%%%%%%%%%%%%%%

% Before a CORE appendix is considered complete, verify ALL items:
%
% □ Answers exactly one fundamental question (WHO/WHAT/HOW/WHEN)
% □ Has complete axiom system with formal statements
% □ Has 80-150+ references organized by category
% □ Has 5-10 critical foundations with steelmanned objections
% □ Has 5-8 open issues with research roadmap
% □ Has bidirectional integration with ALL other CORE appendices
% □ Has comprehensive glossary (20+ symbols)
% □ Has Quick Reference Box
% □ Has at least one detailed worked example
% □ Has Implications for Practice section
% □ Uses CATEGORY-NAME format in title

% =============================================================================
% END OF CORE APPENDIX REQUIREMENTS
% =============================================================================
